%%=============================================================================
%% Voorwoord
%%=============================================================================

\chapter*{\IfLanguageName{dutch}{Woord vooraf}{Preface}}%
\label{ch:voorwoord}

%% TODO:
%% Het voorwoord is het enige deel van de bachelorproef waar je vanuit je
%% eigen standpunt (``ik-vorm'') mag schrijven. Je kan hier bv. motiveren
%% waarom jij het onderwerp wil bespreken.
%% Vergeet ook niet te bedanken wie je geholpen/gesteund/... heeft

%% \lipsum[1-2]

Tijdens een gesprek met ILVO kwam de vraag naar voren of het mogelijk was om een probleem op te lossen
dat zij hadden. Het onderwerp kwam op omdat hun onderzoekers hun berekeningsmodellen willen testen en aanpassen, 
maar ze hebben geen ervaring genoeg met C\# om dit zelf te doen.
De onderzoekers hebben wel ervaring met Python en daarom leek het een goed idee om een Jupyter-notebook 
te maken waarin de onderzoekers kunnen experimenteren met de modellen.

Veel dank aan mijn moeder voor het helpen bij het schrijven van deze bachelorproef,
aan mijn copromoter, Samuel Bosch, en aan mijn promotor, Jan Willem, voor hun begeleiding en feedback tijdens het proces.
